\begin{abstract}
针对异构人工智能任务与能源系统耦合的跨区域调度问题，本文建立算—储—电协同模型。问题一采用固定起点季节基线与 Huber-Ridge 回归预测 GPU 需求，测试期预测总量为 13155.739 等效 GPU，MAE 均值为 25.837 等效 GPU，并以确定性启发式完成 538 个末端任务调度。问题二综合净成本、碳排放、时延和新能源利用率，以任务迁移与开工换时为邻域实施有界局部搜索，完成 50000 个任务的全时域调度；计及新能源外售收入后的净成本为 -421711041.270 元，碳排放为 188.761 tCO2，时延为 5.348 ms，新能源利用率为 67.960\%。问题三建立含充放电、SOC 终端状态及购售电边界的能量流模型，采用购电价格分位规则确定充放电并逐时递推求解；结果表明，成本增加 304223.929 元，绝对爬坡量增加 247.775 MW，既定规则未改善经济性与波动性。问题四建立涵盖成本、碳排放、时延、服务质量、新能源利用率及峰值净购电的多目标模型，采用无量纲加权评分生成 5 个候选，并经约束复核筛选折中方案。
\end{abstract}

\noindent\textbf{关键词：}算电协同；异构任务调度；Huber-Ridge 回归；储能调度；新能源利用；有界局部搜索
